% !TeX root = ../../main.tex
\subsection{疑問詞の種類}
    疑問詞はそもそも次の8つである。

    \begin{definitionbox}{\textbf{疑問詞の種類}}
        \begin{enumerate}
            \item \makebox[7em][l]{\fitblankbf[]{what}}\hbox to 3em{\dotfill\hss}\makebox[11em][l]{\fitblankbf[]{何を}}
            $\Longrightarrow$ 大きく \fitblankbf[]{もの} を問う
            \item \makebox[7em][l]{\fitblankbf[]{who}}\hbox to 3em{\dotfill\hss}\makebox[11em][l]{\fitblankbf[]{誰が}}
            $\Longrightarrow$ \fitblankbf[]{人} を問う
            \item \makebox[7em][l]{\fitblankbf[]{where}}\hbox to 3em{\dotfill\hss}\makebox[11em][l]{\fitblankbf[]{どこで}}
            $\Longrightarrow$ \fitblankbf[]{場所} を問う
            \item \makebox[7em][l]{\fitblankbf[]{when}}\hbox to 3em{\dotfill\hss}\makebox[11em][l]{\fitblankbf[]{いつ}}
            $\Longrightarrow$ \fitblankbf[]{時} を問う
            \item \makebox[7em][l]{\fitblankbf[]{whose}}\hbox to 3em{\dotfill\hss}\makebox[11em][l]{\fitblankbf[]{誰の（もの）}}
            $\Longrightarrow$ \fitblankbf[]{所有者} を問う
            \item \makebox[7em][l]{\fitblankbf[]{which}}\hbox to 3em{\dotfill\hss}\makebox[11em][l]{\fitblankbf[]{どの}}
            $\Longrightarrow$ \fitblankbf[]{選択肢} を問う
            \item \makebox[7em][l]{\fitblankbf[]{why}}\hbox to 3em{\dotfill\hss}\makebox[11em][l]{\fitblankbf[]{なぜ}}
            $\Longrightarrow$ \fitblankbf[]{理由や目的} を問う
            \item \makebox[7em][l]{\fitblankbf[]{how}}\hbox to 3em{\dotfill\hss}\makebox[11em][l]{\fitblankbf[]{どのように}}
            $\Longrightarrow$ \fitblankbf[]{方法や手段} を問う
        \end{enumerate}
    \end{definitionbox}

\subsection{特別な疑問詞の使い方}
    特別な疑問詞の使い方については、大きく分けて次のようなものがある。

    \begin{enumerate}[itemsep=0.8em]
        \item \fitblankbf[]{疑問詞 + 名詞}で1つの疑問詞として扱う使い方
        \item \fitblankbf[]{疑問詞 + 形容詞 [ 副詞 ]}で1つの疑問詞として扱う使い方
        \item 疑問詞が \fitblankbf[]{主語}として使われる使い方
        \item \fitblankbf[]{感嘆文}で使われる使い方
    \end{enumerate}

    それぞれについて、次の小節で深掘りしていくこととする。
\newpage

\subsection{(疑問詞+名詞)で1つの疑問詞として考える}
    (疑問詞+名詞) で1つの疑問詞として扱える疑問詞は\fitblankbf[]{what}、\fitblankbf[]{which}、\fitblankbf[]{whose}の3つのみである。

\subsection{(疑問詞+形容詞 [ 副詞 ])で1つの疑問詞として考える}
    (疑問詞+形容詞 [ 副詞 ]) で1つの疑問詞として扱える疑問詞は\fitblankbf[]{how}のみである。

    具体的には、次のように使う。

\subsection{疑問詞が主語として使われる}
    疑問詞が主語として使われる疑問詞は\fitblankbf[]{what}、\fitblankbf[]{who}の2つのみである。



\vspace{2em}

    以上のパターンを意識して、それらを\fitblankbf[]{1つの疑問詞}として扱い、疑問文を作るということを意識してほしい。
    具体的な疑問文の作り方は、次の節で解説することとする。


\newpage
